
\documentclass{article}
\usepackage{colortbl}
\usepackage{makecell}
\usepackage{multirow}
\usepackage{supertabular}

\begin{document}

\newcounter{utterance}

\centering \large Interaction Transcript for game `taboo', experiment `low\_de', episode 9 with Qwen3.6{-}35B{-}A3B{-}FP8{-}without{-}reasoning.
\vspace{24pt}

{ \footnotesize  \setcounter{utterance}{1}
\setlength{\tabcolsep}{0pt}
\begin{supertabular}{c@{$\;$}|p{.15\linewidth}@{}p{.15\linewidth}p{.15\linewidth}p{.15\linewidth}p{.15\linewidth}p{.15\linewidth}}
    \# & $\;$A & \multicolumn{4}{c}{Game Master} & $\;\:$B\\
    \hline

    \theutterance \stepcounter{utterance}  
    & & \multicolumn{4}{p{0.6\linewidth}}{
        \cellcolor[rgb]{0.9,0.9,0.9}{
            \makecell[{{p{\linewidth}}}]{
                \texttt{\tiny{[P1$\langle$GM]}}
                \texttt{Du spielst ein kooperatives Wortratespiel, in dem du ein Zielwort beschreiben musst, das ein anderer Spieler erraten soll.} \\
\\ 
\texttt{Regeln:} \\
\texttt{(a) Du musst in der Form antworten: HINWEIS: <irgendein Text>. Vermutungen des anderen Spielers beginnen mit VERMUTUNG.} \\
\texttt{(b) Du darfst das Zielwort selbst, Teile davon oder morphologische Varianten davon in deiner Beschreibung nicht verwenden.} \\
\texttt{(c) Zusätzlich gelten die gleichen Regeln für verwandte Wörter, die unten angegeben sind.} \\
\\ 
\texttt{Endbedingungen:} \\
\texttt{(i) Wenn du das Zielwort oder ein verwandtes Wort in deiner Beschreibung verwendest, dann verlierst du.} \\
\texttt{(ii) Wenn der andere Spieler das Zielwort in 3 Versuchen erraten kann, gewinnt ihr beide.} \\
\\ 
\texttt{Lass uns anfangen.} \\
\\ 
\texttt{Dies ist das Zielwort, das du beschreiben musst und das der andere Spieler erraten muss:} \\
\\ 
\texttt{Pharmakologie} \\
\\ 
\texttt{Verwandte Wörter sind:} \\
\\ 
\texttt{{-} arzneimittelkunde} \\
\texttt{{-} adjuvans} \\
\texttt{{-} alchemie} \\
\\ 
\texttt{Wichtig: Du stehst unter Zeitdruck, gib kurze Beschreibungen, die auf den Punkt sind!} \\
            }
        }
    }
    & \\ \\

    \theutterance \stepcounter{utterance}  
    & \multicolumn{4}{p{0.6\linewidth}}{
        \cellcolor[rgb]{0.9,0.9,0.9}{
            \makecell[{{p{\linewidth}}}]{
                \texttt{\tiny{[P1$\rangle$GM]}}
                \texttt{HINWEIS: Dies ist der medizinische Fachbereich, der sich mit der Wirkung, dem Wirkungsmechanismus und der Anwendung von Arzneimitteln im Körper befasst.} \\
            }
        }
    }
    & & \\ \\

    \theutterance \stepcounter{utterance}  
    & & & \multicolumn{2}{p{0.3\linewidth}}{
        \cellcolor[rgb]{0.9,0.9,0.9}{
            \makecell[{{p{\linewidth}}}]{
                \texttt{\tiny{[GM$|$GM]}}
                \texttt{continue} \\
            }
        }
    }
    & & \\ \\

    \theutterance \stepcounter{utterance}  
    & & & \multicolumn{2}{p{0.3\linewidth}}{
        \cellcolor[rgb]{0.9,0.9,0.9}{
            \makecell[{{p{\linewidth}}}]{
                \texttt{\tiny{[GM$|$GM]}}
                \texttt{dies ist der medizinische fachbereich der sich mit der wirkung dem wirkungsmechanismus und der anwendung von arzneimitteln im körper befasst} \\
            }
        }
    }
    & & \\ \\

    \theutterance \stepcounter{utterance}  
    & & \multicolumn{4}{p{0.6\linewidth}}{
        \cellcolor[rgb]{0.9,0.9,0.9}{
            \makecell[{{p{\linewidth}}}]{
                \texttt{\tiny{[GM$\rangle$P2]}}
                \texttt{Du spielst ein kollaboratives Wort{-}Ratespiel, in dem du ein Zielwort erraten musst, das ein anderer Spieler dir beschreibt.} \\
\\ 
\texttt{Du kannst bei jedem Versuch genau einen Tipp abgeben. Du gewinnst, wenn du das Zielwort errätst. Du verlierst, wenn du es in 3 Versuchen nicht erraten kannst.} \\
\\ 
\texttt{Nach jedem Versuch bekommst du einen neuen Hinweis vom anderen Spieler, der mit HINWEIS beginnt.} \\
\\ 
\texttt{Gib deine Tipps ab, indem du einfach das Wort sagst und dabei das folgende Format verwendest: TIPP: <ein Wort>} \\
\\ 
\texttt{Los geht’s.} \\
\\ 
\\ 
\texttt{HINWEIS: Dies ist der medizinische Fachbereich, der sich mit der Wirkung, dem Wirkungsmechanismus und der Anwendung von Arzneimitteln im Körper befasst.} \\
            }
        }
    }
    & \\ \\

    \theutterance \stepcounter{utterance}  
    & & & \multicolumn{4}{p{0.6\linewidth}}{
        \cellcolor[rgb]{0.9,0.9,0.9}{
            \makecell[{{p{\linewidth}}}]{
                \texttt{\tiny{[GM$\langle$P2]}}
                \texttt{TIPP: Pharmakologie} \\
            }
        }
    }
     \\ \\

    \theutterance \stepcounter{utterance}  
    & & & \multicolumn{2}{p{0.3\linewidth}}{
        \cellcolor[rgb]{0.9,0.9,0.9}{
            \makecell[{{p{\linewidth}}}]{
                \texttt{\tiny{[GM$|$GM]}}
                \texttt{continue} \\
            }
        }
    }
    & & \\ \\

    \theutterance \stepcounter{utterance}  
    & & & \multicolumn{2}{p{0.3\linewidth}}{
        \cellcolor[rgb]{0.9,0.9,0.9}{
            \makecell[{{p{\linewidth}}}]{
                \texttt{\tiny{[GM$|$GM]}}
                \texttt{pharmakologie} \\
            }
        }
    }
    & & \\ \\

    \theutterance \stepcounter{utterance}  
    & & & \multicolumn{2}{p{0.3\linewidth}}{
        \cellcolor[rgb]{0.9,0.9,0.9}{
            \makecell[{{p{\linewidth}}}]{
                \texttt{\tiny{[GM$|$GM]}}
                \texttt{end game} \\
            }
        }
    }
    & & \\ \\

\end{supertabular}
}

\end{document}
